\documentclass[11pt]{article}
\input{../../notes-en/structure}
\hypersetup{pdftitle={},pdfauthor={}}
\begin{document}
\begin{center}
{\LARGE\sffamily\bfseries\color{deepblue} Week 5 Bonus Quiz}\\[.2cm]
{\large Functions that work on lists, without \texttt{LinearAlgebra}}
\end{center}

\noindent Both questions are optional bonus questions. The grading environment has
no \texttt{LinearAlgebra}, so \texttt{dot} and \texttt{norm} are unavailable: write
the loop yourself.

\section*{Question 2 --- normalized inner product}
Write a function that receives two lists of numbers and returns their inner
product, normalized by their lengths:
\[
\frac{a\cdot b}{|a|\,|b|}.
\]
Signature:
\begin{lstlisting}
function norm_inner_product(v1, v2)
    # ...
end
\end{lstlisting}
Return value: a number. You may assume the input is two lists containing only
\texttt{Int} or \texttt{Float64} values.
\begin{itemize}
\item If the lists have different lengths, return \texttt{nothing}.
\item If the lists are empty, return \texttt{0}.
\end{itemize}
A zero vector makes the normalized product undefined; the tests do not include
that case.

Example: for \texttt{v1 = [3.0, 4.0]} and \texttt{v2 = [6.0, 8.0]} the result is
\texttt{1.0}. For \texttt{[1.0, 0.0]} and \texttt{[0.0, 1.0]} it is \texttt{0.0}.

\section*{Question 4 --- counting orthogonal pairs}
Two vectors are orthogonal when their inner product is zero. Write a function
that receives a list of vectors (a list of lists) and returns the number of pairs
of vectors that are orthogonal to each other:
\begin{lstlisting}
function num_of_orthogonal(vectors)
    # ...
end
\end{lstlisting}
Return value: the number of orthogonal pairs in the list.

\begin{lstlisting}
num_of_orthogonal([[1, 0, 0], [0, 1, 0], [0, 0, 1]])              # 3
num_of_orthogonal([[0, 0, 0], [0, 1, 0], [0, 0, 1]])              # 3
num_of_orthogonal([[0, 0], [1, 2], [10, 5]])                      # 2
num_of_orthogonal([[1, 1, 1, 1], [2, 1, 3, 3], [0, 0, 100, 33],
                   [8, 8, 8, 1.5], [9, 9, 9, 9]])                 # 0
\end{lstlisting}
You may assume that the input is a list, that the inner lists have equal length,
and that they contain only numbers. Count each pair once: the pair $(i,j)$ and
the pair $(j,i)$ are the same pair, and a vector is not paired with itself.
Avoid duplicated code; use one inner-product helper for both questions.

\begin{center}
\begin{tabular}{lc}
\toprule
Call & Result \\
\midrule
\texttt{norm\_inner\_product([1.0, 0.0], [-1.0, 0.0])} & \texttt{-1.0} \\
\texttt{norm\_inner\_product([1, 2, 3], [1, 2])} & \texttt{nothing} \\
\texttt{norm\_inner\_product(Float64[], Float64[])} & \texttt{0} \\
\texttt{num\_of\_orthogonal} on the 4D basis & \texttt{6} \\
\bottomrule
\end{tabular}
\end{center}
\end{document}
